\section{Related work and positioning}

\subsection{Controlled smoothness in finite-difference PLS}

The direct methodological foundation is the controlled-smoothness formulation of finite-difference PLS. Guerrero \cite{guerrero2007penalized} derives a statistical representation of the penalized trend estimator and introduces a trace-based smoothness index linking the penalty parameter to an interpretable amount of smoothness. Guerrero \cite{guerrero2008smoothness} develops this into a user-controlled trend-estimation procedure in which a desired smoothness level is chosen first and the corresponding penalty is recovered afterward.

Cort\'es-Toto et al.\ \cite{cortes2017smoothness} study what smoothness is produced when the PLS penalty is selected by CV, GCV, AICc, and BIC. Their results make an important distinction for the present paper: a rule for selecting $\lambda$ and the amount of smoothness ultimately induced by that rule are not the same object.

The present paper reverses the controlled-smoothness direction:
\[
\text{user chooses }S\longrightarrow\lambda(S)
\]
becomes
\[
\text{chronological forecast loss}\longrightarrow S^\star.
\]
The percentage of smoothness remains the interpretable object, but it is selected endogenously for the forecast task.

\subsection{Predictive selection of smoothing and tuning parameters}

Forecast-based tuning is established in other smoother families. Hart \cite{hart1994tscv} considers kernel smoothing with serially correlated errors and proposes time-series cross-validation (TSCV), selecting the bandwidth through one-step-ahead predictive performance. This is an important conceptual precedent because it distinguishes good trend recovery from good prediction under dependence.

The object considered here is different. Hart selects a kernel bandwidth and explicitly models serial dependence in the errors. We select the smoothness of a finite-difference PLS trend, retain the model's native continuation rule, and allow an arbitrary $h$-step future block to define the tuning target. The connection to Hart is therefore the predictive principle rather than estimator equivalence.

Forecast-error tuning is also standard in exponential smoothing. Taylor \cite{taylor2004stes} describes data-driven smoothing-parameter estimation through ex post one-step-ahead forecast errors and develops an adaptive smooth-transition formulation. More recent \emph{Journal of Forecasting} papers apply the same predictive principle to other tuning problems. Wolff and Echterling \cite{wolff2024stock} choose regularization and machine-learning tuning parameters using validation forecast errors. Franjic and Schweikert \cite{franjic2025predictor} propose a cross-validation strategy that connects predictor preselection directly to nowcast performance. Xu et al.\ \cite{xu2025vix} compare bandwidth-selection schemes in a time-varying-parameter HAR model and show that bandwidth choice affects forecast performance.

These papers establish that forecast-oriented tuning is a legitimate methodological target. The narrower gap addressed here is the selection of a controlled smoothness percentage inside the finite-difference PLS family.

\subsection{Trend filtering for forecasting}

Zafar et al.\ \cite{zafar2022filter} explicitly evaluate trend-filtering methods by the short-term forecasts they enable. Their study reinforces the idea that a trend representation should be judged partly by predictive consequences rather than only by visual smoothness or in-sample fit. Our objective is complementary: rather than introduce a new trend representation, we keep finite-difference PLS fixed and optimize the amount of smoothness within that family.

\subsection{Chronological forecast evaluation}

The proposed criterion uses pseudo-out-of-sample loss as a tuning objective. Tashman \cite{tashman2000oos} reviews out-of-sample forecast evaluation schemes, while Bergmeir and Ben\'itez \cite{bergmeir2012cv} discuss cross-validation for time-series predictors. Stan\v{e}k \cite{stanek2023oos} formalizes out-of-sample loss estimation and distinguishes rolling and fixed evaluation schemes. His formulation also makes explicit that the relevant forecast loss depends on the forecast horizon, model memory, and updating design.

This is closely aligned with $F_{d,L,h}(S)$. Smoothness is selected for a specified horizon and update protocol, and the information set available at each origin is part of the estimator's definition.
